\section{战略专家二：AGI 时代的终端与下一代苹果}

本章讨论访谈中最宏大的战略想象：理想是否能成为 AGI 时代的苹果。李想把 AGI 时代终端定义为具备四类能力：360 度感知物理世界、认知决策、Action、反思反馈。这个定义把“终端”从屏幕设备扩展到物理世界入口：车可以是终端，机器人也可以是终端，家庭设备和工厂设备也可能成为终端。

\begin{figure}[H]
\centering
\includegraphics[width=0.96\textwidth]{agi-era-apple-route.png}
\caption{AGI 时代的苹果路线：车、机器人、操作系统和模型能力共同构成入口。自制概念图，依据 01:38:20--01:46:40 与 02:09:16--02:09:21 对谈内容整理。}
\end{figure}

\begin{knowledgebox}{读图：车不是唯一终点，而是物理入口}
图中的车、VLA、机器人、操作系统和 AGI Apple 不是线性口号。车是现有最大规模物理入口，VLA 是驾驶大脑，机器人扩展到家庭和工厂，操作系统连接软件和工具，最终才可能形成下一代平台。
\end{knowledgebox}

\subsection{软件能力：操作系统、虚拟机和工具}

上一节把车和机器人放进下一代平台想象，本节先拆软件底座。李想认为，AGI 时代终端需要不同的软件能力。第一，要有操作系统能力，能让 AI 在物理和数字世界中稳定运行。第二，要有虚拟机或执行环境能力，让 Agent 在本地电脑、云端和设备之间运行。第三，要有工具生态，因为 Agent 需要调用工具，而工具本身也会为 AI 重新设计。

\begin{importantbox}{为什么操作系统重新变重要}
当 AI 只是应用时，操作系统像背景；当 AI 要持续观察、调用工具、控制设备、记录状态并执行动作时，操作系统就变成安全、性能、权限和生态的核心。
\end{importantbox}

\subsection{硬件能力：感知、算力和本体}

软件底座解决“怎样运行”，本节转向“运行在什么身体上”。AGI 终端也需要硬件能力。车有传感器、算力、能源、空间和移动能力；机器人有本体、执行器、触觉和工厂/家庭场景；家庭设备可能需要统一感知和统一大脑。访谈里李想没有把机器人路线说成唯一答案，而是承认两种可能：人形机器人可以用人类工具，也可以改造工具和环境，让非人形设备完成任务。

\begin{knowledgebox}{机器人路线的两种想象}
\begin{tabularx}{\textwidth}{p{0.25\textwidth}p{0.34\textwidth}X}
\toprule
路线 & 优势 & 难点 \\
\midrule
人形机器人 & 适配人类现有工具和空间 & 本体复杂、成本高、控制难。 \\
改造环境/工具 & 更容易在单任务中做到高可靠 & 需要重新设计家庭、工厂或设备生态。 \\
统一大脑 & 可复用模型、记忆和任务理解 & 不同本体之间动作空间差异大。 \\
\bottomrule
\end{tabularx}
\end{knowledgebox}

\subsection{销量与 AI 价值的错位}

市场现在仍主要用销量看理想，因为 AI 能力带来的云服务价值、软件价值和平台价值还没有充分表现出来。这个错位并不罕见：在平台转型早期，外部常常先用旧指标评价新能力。真正的拐点会出现在 AI 能力能改善体验、降低成本、带来收入或打开新终端形态时。

\begin{warningbox}{“成为苹果”不是品牌比喻，而是平台比喻}
如果只把苹果理解为高端品牌，就会误读这段讨论。李想谈的是平台能力：硬件、操作系统、开发生态、用户入口、工具链和商业模式。车企要成为 AGI 时代的苹果，必须证明自己不只是卖车。
\end{warningbox}

\subsection{本章小结}

AGI 时代终端的想象把理想的战略从汽车扩展到平台。车是入口，VLA 是能力，Agent OS 是软件层，机器人是延伸，本质问题是公司能否把物理世界入口变成智能平台。

